\documentclass[11pt]{letter}

\usepackage[margin=0.75in]{geometry}
\usepackage{setspace}
\usepackage{hyperref}
\usepackage{xcolor}
\usepackage{enumitem}
\usepackage{parskip}

\hypersetup{
    colorlinks=true,
    urlcolor=black,
    linkcolor=black
}

\setlength{\parindent}{0pt}
\setstretch{1.05}
\date{}
\begin{document}

\begin{letter}{}

\opening{Dear Hiring Committee,}

I saw the job post on the Sola careers page and thought I'd be a good fit for it.

I am a Master's student in Computer Science at the University of Cincinnati, graduating December 2026, looking for a \textbf{Software Engineering} role at \textbf{Sola}. With production experience spanning distributed backend infrastructure and hands-on AI research, I am excited about the opportunity to contribute to \textbf{Sola} and help build reliable, observable execution on top of non-deterministic AI models.

\textbf{Here are reasons that make me a great fit for the role:}

\begin{itemize}[leftmargin=*, itemsep=0.45em]
    \item \textbf{Owning AI features end-to-end, from experimentation to production:} My graduate research designs and evaluates a two-stage memory-augmented LLM pipeline for Android automation, taking it from an early prototype through a structured evaluation against a published baseline (GPT-4o + GroundingDINO + CLIP), achieving a 25.7x latency reduction while preserving task fidelity across 20+ real-world apps.

    \item \textbf{Building and operating reliable distributed backend systems:} At Siemens Industry Software, I migrated 40+ production environments to Kubernetes using ArgoCD, Helm, and GitLab CI/CD, and designed an Ansible based Blue Green deployment framework that cut upgrade downtime 50\% with automated rollback within 15 minutes.

    \item \textbf{Proficiency in Python and modern ML frameworks:} I am proficient in Python and TypeScript, with hands-on experience across PyTorch, scikit-learn, and Vertex AI, including building and benchmarking a PyTorch LSTM classifier on 400K+ headlines and provider-agnostic LLM evaluation harnesses spanning Gemini, GPT-4o, Llama, and Qwen.

    \item \textbf{Building observability into production systems:} At Siemens, I diagnosed live production issues across 40+ environments using CloudWatch and Datadog log correlation, and built a self-service developer platform adopted by 15 engineers, the same instinct for monitoring and measuring system health I would bring to evals and observability infrastructure for AI systems in production.

    \item \textbf{Comfort with ambiguity at an early stage:} I have shipped systems end to end in production, from Terraform-based AWS provisioning to a self-service CLI and dashboard, and enjoy the ambiguity of figuring out what is feasible to build rather than working from a fixed spec. A small, in-person team turning cutting-edge AI research into shipped features is exactly the environment I want to be in.
\end{itemize}

I am available to work with you after my graduation in December 2026, and am glad to relocate to New York City. As a Master's student at the University of Cincinnati and someone who enjoys turning AI research into reliable production systems, I am enthusiastic about contributing to Sola and helping bring state-of-the-art agentic capabilities to production at scale.

Please find my CV attached below. I look forward to hearing from you. Sola's mission of building reliable, deterministic execution on top of non-deterministic AI models is exactly the kind of hard, high-impact problem I want to help solve.

{\setlength{\parskip}{0pt}
Regards,\\
Tanmay Bhuskute,\\
bhuskutv@mail.uc.edu $\mid$ tanmay.v.bhuskute@gmail.com $\mid$ +1 (732) 858-3027,\\
3424 Brookline Ave, Flat 19, Cincinnati, OH 45220.
}

\end{letter}

\end{document}
